%%=============================================================================
%% Methodologie
%%=============================================================================

\chapter{\IfLanguageName{dutch}{Methodologie}{Methodology}}%
\label{ch:methodologie}

%% TODO: In dit hoofstuk geef je een korte toelichting over hoe je te werk bent
%% gegaan. Verdeel je onderzoek in grote fasen, en licht in elke fase toe wat
%% de doelstelling was, welke deliverables daar uit gekomen zijn, en welke
%% onderzoeksmethoden je daarbij toegepast hebt. Verantwoord waarom je
%% op deze manier te werk gegaan bent.
%% 
%% Voorbeelden van zulke fasen zijn: literatuurstudie, opstellen van een
%% requirements-analyse, opstellen long-list (bij vergelijkende studie),
%% selectie van geschikte tools (bij vergelijkende studie, "short-list"),
%% opzetten testopstelling/PoC, uitvoeren testen en verzamelen
%% van resultaten, analyse van resultaten, ...
%%
%% !!!!! LET OP !!!!!
%%
%% Het is uitdrukkelijk NIET de bedoeling dat je het grootste deel van de corpus
%% van je graduaatsproef in dit hoofstuk verwerkt! Dit hoofdstuk is eerder een
%% kort overzicht van je plan van aanpak.
%%
%% Maak voor elke fase (behalve het literatuuronderzoek) een NIEUW HOOFDSTUK aan
%% en geef het een gepaste titel.

Dit hoofdstuk beschrijft de aanpak op hoofdlijnen. Het oorspronkelijke plan vertrok van basis naar uitbreiding: requirements, databankontwerp, backend, authenticatie, ticketbeheer, berichten, frontend, adminschermen, uploads, e-mail, testen en documentatie. Tijdens de uitvoering werd dit bijgestuurd op basis van mentorfeedback, technische keuzes en onderhoudbaarheid.

\section{Aanpak in fasen}

\begin{itemize}[itemsep=0.25\baselineskip, topsep=0.25\baselineskip, leftmargin=*]
  \item \textbf{Scope en analyse.} Eerst werden rollen, gegevens en workflows bepaald vanuit klant, administrator en systeem. Dit leverde de basis op voor gebruikers, tickets, berichten, bijlagen en de statusflow.
  \item \textbf{Technische basis.} Daarna werd het proof of concept opgebouwd met ASP.NET Core, Entity Framework Core en MySQL. De backend kreeg meteen lagen, DTO's, repositories, services, migraties, JWT-authenticatie, foutafhandeling en het Options Pattern~\autocite{MicrosoftOptionsPattern2026}.
  \item \textbf{Frontend en feedback.} De frontend werd gebouwd met Expo~\autocite{ExpoDocs2026}. Formik en Yup ondersteunden formulieren en validatie~\autocite{FormikDocs2026,YupDocs2026}. Mentorfeedback leidde tot chatbubbels, afbeeldingspreview, uploadlimiet, responsieve lay-out, browsernavigatie, zoekvelden, filters en paginatie.
  \item \textbf{Controle en documentatie.} Ten slotte werd realistische testdata toegevoegd voor meerdere klanten, tickets en langere gesprekken. Filters, statusgroepen, paginatie, databankopstart, codecommentaar en README-bestanden werden gecontroleerd, zodat de applicatie makkelijker te onderhouden blijft.
\end{itemize}
